% !TeX program = xelatex
% !BIB program = biber
% !TEX encoding = UTF-8 Unicode
% !TeX spellcheck = de_DE

% Template von Leon Willens adaptiert von Michael Rademacher
% Article-Class mit 11pt Schriftgröße, Zielformat DIN A4 und Titelseite
\documentclass[11pt,a4paper,titlepage,final,table]{article}

% Meta-Commands
% !TEX root = "./main.tex"
% Meta-Daten, die mehrfach verwendet werden können
% Angaben Autor*in
\newcommand{\AuthorFirstName}{Max}
\newcommand{\AuthorLastName}{Mustermann}
\newcommand{\AuthorStudentNumber}{1234567}
\newcommand{\AuthorStreet}{Musterstraße 1}
\newcommand{\AuthorZipCode}{12345}
\newcommand{\AuthorLocality}{Musterstadt}
\newcommand{\AuthorEmail}{max.mustermann@smail.inf.h-brs.de}
% Kombinationen
\newcommand{\Author}{\AuthorFirstName\ \AuthorLastName}
\newcommand{\AuthorAddress}{\AuthorStreet \newline \AuthorZipCode\ \AuthorLocality}
% Für externe Arbeiten
\newcommand{\ExternalCompany}{Musterfirma}
\newcommand{\Location}{Musterstadt}
% Angaben zum Studium
\newcommand{\CourseOfStudies}{Musterstudium}
\newcommand{\CourseOfStudiesDegree}{Bachelor / Master}
% Angaben zum Dokument
\newcommand{\DocumentType}{Thesis}
\newcommand{\DocumentTitle}{Mustertitel}
\newcommand{\DocumentSubject}{\DocumentType im \CourseOfStudiesDegree\-Studiengang \CourseOfStudies}
\newcommand{\Keywords}{}
\newcommand{\FirstSupervisor}{Prof. Dr. Martina Mustermann}
\newcommand{\SecondSupervisor}{Prof. Dr. Maria Mustermann}
% Entweder festes Datum oder \today
\newcommand{\DateOfSubmission}{\today}
% Eigentlich obsolet, aber vielleicht kann man das ja auch für andere Fachbereiche benutzen.
\newcommand{\Department}{Computer Science}
\newcommand{\DepartmentGER}{Informatik}  
% STY-Datei mit Package- und Command-Definitionen   		
\usepackage{./sty/thesis-de}

% Hier Literatur-Ressourcen einfügen
\addbibresource{./lit/lit.bib}
% \addbibresource{./lit/foo.bib}

\begin{document}
\maketitle
% Eidesstattliche Erklärung - Kommentieren falls nicht benötigt
\input{./metapages/affidavit}

\clearpage
\pagenumbering{roman}
\setcounter{page}{2}
\pagestyle{plain}
\parskip 0.5em

% Inhaltsverzeichnis
\tableofcontents
\clearpage

% Abbildungsverzeichnis
\listoffigures
\clearpage

% Tabellenverzeichnis
\listoftables
\clearpage

\parskip 1em

% Abstract / Kurzfassung - Kommentieren falls nicht benötigt
\input{./section/abstract}

\clearpage
\pagenumbering{arabic}

\lstset{
	basicstyle=\ttfamily\small,
	keywordstyle=\bfseries,
	language=[Sharp]C
}

% Hier Inhalte einfügen
% Einleitung
\section{Einleitung}
\label{sec:intro}
Dieses \LaTeX-Template wurde konzipiert für die Verfassung von Abschlussarbeiten im Fachbereich Informatik der \ac{H-BRS}.
Es orientiert sich hinsichtlich der formalen Vorgaben an den Regeln für wissenschaftliches Arbeiten~\cite{Kees2017}.

% Abschnitt 2
%\input{./section/foo}

% Literatur
\newgeometry{top=25mm,right=20mm,bottom=20mm,left=20mm}
\printbibliography[heading=bibintoc]

\section*{Abkürzungsverzeichnis}
\begin{acronym}[]
%\addcontentsline{toc}{section}{List of Abbreviations}
%\section*{List of abbreviations}
\setlength{\itemsep}{-\parsep}


\acro{H-BRS}{Hochschule Bonn-Rhein-Sieg}
\end{acronym}



% Anhang
\appendix
% Für Deutsch
\renewcommand{\appendixname}{Anhang}
\renewcommand{\appendixtocname}{Anhang}
\renewcommand{\appendixpagename}{Anhang}
\begin{appendices}
% Anhänge hier
\end{appendices}

\end{document}